% 3.8 例子：立方密堆积与金刚石结构
\section{例子：立方密堆积与金刚石结构}

本节给出两个例子来说明前几节的理论；一个是点式空间群——立方密堆积 $Fm3m\ (O^{5}_{h})$，另一个是非点式空间群——金刚石结构的空间群 $Fd3m\ (O^{7}_{h})$。这两个空间群分别最先由 Bouckaert, Smoluchowski, and Wigner (1936) 与 Herring (1942) 处理过。

二者都建立在面心立方（$F$）布拉维格子 $\Gamma^{f}_{c}$ 上，故由表 3.1，两种情形的格子平移都是
\begin{equation*}
\left\{
\begin{array}{l}
\mathbf{t}_{1} = \frac{1}{2}(0, a, a),\\
\mathbf{t}_{2} = \frac{1}{2}(a, 0, a),\\
\mathbf{t}_{3} = \frac{1}{2}(a, a, 0).
\end{array}
\right.
\tag{3.8.1}
\end{equation*}
两种情形的同形点群 $\mathbf{F}$ 都是 $m3m\ (O_{h})$，即完整立方群。

由表 3.3 可见倒格子矢量为
\begin{equation*}
\left\{
\begin{array}{l}
\mathbf{g}_{1} = (2\pi/a)(-1, 1, 1),\\
\mathbf{g}_{2} = (2\pi/a)(1, -1, 1),\\
\mathbf{g}_{3} = (2\pi/a)(1, 1, -1),
\end{array}
\right.
\tag{3.8.2}
\end{equation*}
布里渊区示于图 3.14。

由表 3.7 可见，对 $Fm3m\ (O^{5}_{h})$，$m3m\ (O_{h})$ 的每个转动算符都只与纯格子平移相联系，于是关于平移群 $\mathbf{T}$ 的一个合适的左陪集分解为
\begin{equation*}
Fm3m\left(O^{5}_{h}\right) = \sum_{R}\{R \mid 000\}\mathbf{T}
\tag{3.8.3}
\end{equation*}
其中对 $R$ 的求和取遍 $m3m\ (O_{h})$ 的所有元素。而对 $Fd3m\ (O^{7}_{h})$，只有属于子群 $\bar{4}3m\ (T_{d})$ 的算符与纯格子平移相联系，相应的左陪集分解为
\begin{equation*}
Fd3m\left(O^{7}_{h}\right) = \sum_{R}\{R \mid 000\}\mathbf{T} + \sum_{R}\{RI \mid {\textstyle\frac{1}{4}}{\textstyle\frac{1}{4}}{\textstyle\frac{1}{4}}\}\mathbf{T}
\tag{3.8.4}
\end{equation*}
此时对 $R$ 的求和取遍 $\bar{4}3m\ (T_{d})$ 的所有元素，且如通常，$I$ 是反演。与不在 $\bar{4}3m\ (T_{d})$ 中的 $m3m\ (O_{h})$ 元素相对应的那些陪集，其陪集代表总可取为 $\mathbf{v} = \frac{1}{4}(\mathbf{t}_{1} + \mathbf{t}_{2} + \mathbf{t}_{3}) = \frac{1}{4}(a, a, a)$（在 $Oxyz$ 轴下）。算符 $R$ 作用在（$Oxyz$ 轴下的）非零平移 $(p, q, r)$ 上的效果可由表 1.4 查出，例如 $C^{-}_{4z}(p, q, r) = (q, -p, r)$。算符 $R$ 作用在布拉维格子基本矢量 $\mathbf{t}_{1}$、$\mathbf{t}_{2}$、$\mathbf{t}_{3}$ 上的效果可由表 3.2 查出，例如对立方 $F$ 格子 $C^{-}_{4z}\mathbf{t}_{1} = \mathbf{t}_{2}$、$C^{-}_{4z}\mathbf{t}_{2} = \mathbf{t}_{2} - \mathbf{t}_{3}$、$C^{-}_{4z}\mathbf{t}_{3} = -\mathbf{t}_{1} + \mathbf{t}_{2}$。

在表 3.11 中我们列出了布里渊区基本域中我们求 $Fm3m\ (O^{5}_{h})$ 与 $Fd3m\ (O^{7}_{h})$ 的表示时所需要的对称点、对称线和对称面以及一个一般点。它们就是 3.7 节的点 $\mathbf{k}_{1}$。两种情形下，由于 $\mathbf{F}$ 是布拉维格子的全对称点群，表示域与基本域重合，因此就是图 3.14 所示的区域。

\begin{figure}[H]
\centering
\includegraphics[width=0.88\textwidth]{c3_t311.png}
\caption*{\textbf{表 3.11}\ $\Gamma^{f}_{c}$ 的布里渊区的基本域（原书扫描占位）。}
\end{figure}

表 3.11 的注（译文）：

\noindent（i）第 1、2、3 列在表的每一行中分别列出点 $\mathbf{k}_{1}$、它相对于轴 $\Gamma\mathbf{g}_{1}\mathbf{g}_{2}\mathbf{g}_{3}$ 的坐标，以及它以 $2\pi/a$ 为单位、相对于轴 $\Gamma k_{x}k_{y}k_{z}$ 的坐标。点 $\mathbf{k}_{1}$ 选为表示的星的生成元。

\noindent（ii）第 4 列在对应于 $\mathbf{k}_{1}$ 的行中列出 $\mathbf{k}_{1}$ 的星中元素的数目 $q$，以及字母 I 或 S——视 $\mathbf{k}_{1}$ 是布里渊区的内部点还是表面点而定。

\noindent（iii）第 5 列在对应于 $\mathbf{k}_{1}$ 的行中列出小陪群 $\overline{\mathbf{G}}^{\mathbf{k}_{1}}$ 的 $b$ 个元素。所有情形都有 $bq = 48$，即 $m3m\ (O_{h})$ 的阶。

\noindent（iv）本表对两个空间群乃至一切基于 $\Gamma^{f}_{c}$ 且同形点群为 $m3m\ (O_{h})$ 的空间群都是共同的；另见表 3.6 的相应部分。表 3.11 中有而表 3.6 中没有的点，是对称面上的点与一般点。除本例外我们将不再列出对称面，因为如我们将看到的，它们很容易处理。

现在考虑把小群 $\mathbf{G}^{\mathbf{k}_{1}}$ 关于平移群 $\mathbf{T}$ 作左陪集分解。对 $Fm3m$ 有
\begin{equation*}
\mathbf{G}^{\mathbf{k}_{1}} = \sum_{S}\{S \mid 000\}\mathbf{T}
\tag{3.8.5}
\end{equation*}
其中对 $S$ 的求和取遍 $\overline{\mathbf{G}}^{\mathbf{k}_{1}}$ 的所有元素（这些群见表 3.11）。而对 $Fd3m$ 有
\begin{equation*}
\mathbf{G}^{\mathbf{k}_{1}} = \sum_{P}\{P \mid 000\}\mathbf{T} + \sum_{Q}\{Q \mid {\textstyle\frac{1}{4}}{\textstyle\frac{1}{4}}{\textstyle\frac{1}{4}}\}\mathbf{T},
\tag{3.8.6}
\end{equation*}
其中对 $P$ 的求和取遍 $\overline{\mathbf{G}}^{\mathbf{k}_{1}}$ 与 $\bar{4}3m\ (T_{d})$ 共有的元素，对 $Q$ 的求和取遍 $\overline{\mathbf{G}}^{\mathbf{k}_{1}}$ 的其余元素。

\subsection*{Fm3m}

$Fm3m$ 可以很快处理完。给定任一点 $\mathbf{k}_{1}$，需要确定 $\overline{\mathbf{G}}^{\mathbf{k}_{1}}$ 带因子系统 $\exp(-\mathrm{i}\mathbf{g}_{i} \cdot \mathbf{w}_{j})$ 的射影表示，其中 $\mathbf{g}_{i}$ 由式 (3.7.11) 给出，$\mathbf{w}_{j}$ 是分解 (3.8.5) 中与 $S_{j}$ 相联系的平移。立即看出对所有 $j$ 有 $\mathbf{w}_{j} = 0$。因此因子系统完全由单位元组成，属于 $\overline{\mathbf{G}}^{\mathbf{k}_{1}}$ 的乘子的恒等元所对应的类。于是合适的表示 $\mathbf{D}^{\mathbf{k}_{1}}_{p}$ 就是 $\overline{\mathbf{G}}^{\mathbf{k}_{1}}$ 的普通矢量表示；而且由于 $\overline{\mathbf{G}}^{\mathbf{k}_{1}}$ 是点群（诚然某些情形下定向非标准），可以从表 2.2 得到 $\mathbf{D}^{\mathbf{k}_{1}}_{p}$。

$\mathbf{G}^{\mathbf{k}_{1}}$ 的小表示 $\Gamma^{\mathbf{k}_{1}}_{p}$ 现在由式 (3.7.9) 立即得到。$\Gamma^{\mathbf{k}_{1}}_{p}$ 的矩阵在所有情形都与 $\mathbf{D}^{\mathbf{k}_{1}}_{p}$ 的矩阵通过方程
\begin{equation*}
\Gamma^{\mathbf{k}_{1}}_{p}(\{S \mid \mathbf{t}\}) = \exp(-\mathrm{i}\mathbf{k}_{1} \cdot \mathbf{t})\mathbf{D}^{\mathbf{k}_{1}}_{p}(S)
\tag{3.8.7}
\end{equation*}
相联系。因此只需指明点群 $\overline{\mathbf{G}}^{\mathbf{k}_{1}}$ 即可把 $Fm3m$ 处理完毕。结果如下：对 $\Gamma$，$m3m\ (O_{h})$；对 $X$，$4/mmm\ (D_{4h})$；对 $L$，$\bar{3}m\ (D_{3d})$；对 $W$，$\bar{4}2m\ (D_{2d})$；对 $\Delta$，$4mm\ (C_{4v})$；对 $\Lambda$，$3m\ (C_{3v})$；对 $\Sigma$，$mm2\ (C_{2v})$；对 $S$，$mm2\ (C_{2v})$；对 $Z$，$mm2\ (C_{2v})$；对 $Q$，$2\ (C_{2})$；对 $C$，$m\ (C_{1h})$；对 $O$，$m\ (C_{1h})$；对 $J$，$m\ (C_{1h})$；对 $B$，$m\ (C_{1h})$；对 $A$，$1\ (C_{1})$。

\subsection*{Fd3m}

金刚石的情形就不那么容易了。首先注意到，在分解 (3.8.6) 中，若陪集代表的转动部分 $P$ 是元素 $E$、$C^{\pm}_{3j}$、$C_{2m}$、$\sigma_{dp}$ 或 $S^{\pm}_{4m}$ 之一，则其平移部分 $\mathbf{w}_{j} = 0$；而若其转动部分 $Q$ 是元素 $I$、$S^{\mp}_{6j}$、$\sigma_{m}$、$C_{2p}$ 或 $C^{\pm}_{4m}$ 之一，则其平移部分 $\mathbf{w}_{j} = \mathbf{v} = \frac{1}{4}(a, a, a)$。因此因子系统 $\exp(-\mathrm{i}\mathbf{g}_{i} \cdot \mathbf{w}_{j})$ 在下述两种情形完全由单位元组成：对所有 $i$ 有 $\mathbf{g}_{i} = \mathbf{0}$；或者小陪群 $\overline{\mathbf{G}}^{\mathbf{k}_{1}}$ 的元素全部形如 $P$ 而没有形如 $Q$ 的元素。前者在 $\mathbf{k}_{1}$ 是布里渊区内点时成立（$\Gamma, \Delta, \Lambda, \Sigma, C, O, J, A(\mathrm{I})$），后者在 $\overline{\mathbf{G}}^{\mathbf{k}_{1}}$ 完全由 $\bar{4}3m\ (T_{d})$ 中的元素组成时成立（$\Lambda, C, J, A(\mathrm{S})$）。这些情形所需的射影表示 $\mathbf{D}^{\mathbf{k}_{1}}_{p}$ 就是点群 $\overline{\mathbf{G}}^{\mathbf{k}_{1}}$ 的普通表示。

于是，在 $Fd3m$ 中若 $\mathbf{k}_{1} = \Gamma, \Delta, \Lambda, \Sigma, C, O, J$ 或 $A$，求 $\mathbf{D}^{\mathbf{k}_{1}}_{p}$ 时只需在表 2.2 中查 $\overline{\mathbf{G}}^{\mathbf{k}_{1}}$ 的表示即可。与这些点对应的点群 $\overline{\mathbf{G}}^{\mathbf{k}_{1}}$ 的名称已在上面处理 $Fm3m$ 时给出。然而，即使对这些点，$Fd3m$ 与 $Fm3m$ 之间也有本质差别：得到 $\mathbf{D}^{\mathbf{k}_{1}}_{p}$ 之后，小表示 $\Gamma^{\mathbf{k}_{1}}_{p}$ 的导出方式与 $Fm3m$ 不完全一样。此时由式 (3.7.9) 可见
\begin{equation*}
\Gamma^{\mathbf{k}_{1}}_{p}(\{P \mid \mathbf{t}\}) = \exp(-\mathrm{i}\mathbf{k}_{1} \cdot \mathbf{t})\mathbf{D}^{\mathbf{k}_{1}}_{p}(P),
\tag{3.8.8}
\end{equation*}
但
\begin{equation*}
\Gamma^{\mathbf{k}_{1}}_{p}(\{Q \mid \mathbf{t} + \mathbf{v}\}) = \exp\left(-\mathrm{i}\mathbf{k}_{1} \cdot (\mathbf{t} + \mathbf{v})\right)\mathbf{D}^{\mathbf{k}_{1}}_{p}(Q),
\tag{3.8.9}
\end{equation*}
其中 $P$ 与 $Q$ 的含义与分解 (3.8.6) 中相同。由于习惯上列出的是小表示 $\Gamma^{\mathbf{k}_{1}}_{p}$ 而非 $\mathbf{D}^{\mathbf{k}_{1}}_{p}$，表中出现的某些元素带有一个额外因子 $\exp(-\mathrm{i}\mathbf{k}_{1} \cdot \mathbf{v})$，它不出现在相应点群的特征标表中。（值得指出，对 $\Gamma$ 这个因子总是 1，因此属于 $\Gamma$ 的非点式空间群的表示与相应点式空间群的表示相同。）

现在详细考虑布里渊区的每个表面点 $B, Q, Z, S, L, W$ 和 $X$。我们将不加解释地使用 3.7 节后面各段的记号。

\subsection*{B}

$\mathbf{G}^{B}$ 的陪集代表：$\{E \mid 000\}$，$\{\sigma_{y} \mid {\frac{1}{4}}{\frac{1}{4}}{\frac{1}{4}}\}$。

\begin{center}
$EB = B$，于是 $\mathbf{g}_{E} = 0$。

$\sigma_{y}B = B - \mathbf{g}_{1} - \mathbf{g}_{3}$，于是 $\mathbf{g}_{\sigma_{y}} = -\mathbf{g}_{1} - \mathbf{g}_{3} = (-1, 0, -1)$。
\end{center}

因子系统：$\mu(E, E) = 1$，$\mu(E, \sigma_{y}) = 1$，$\mu(\sigma_{y}, E) = 1$，以及 $\mu(\sigma_{y}, \sigma_{y}) = \exp\left(2\pi\mathrm{i}\left(\frac{1}{4} + 0 + \frac{1}{4}\right)\right) = -1$。

中央延拓群 $\overline{\mathbf{G}}^{B*}$ 由四个元素 $(E, 0)$、$(E, 1)$、$(\sigma_{y}, 0)$、$(\sigma_{y}, 1)$ 组成，乘法表为
\begin{center}
\begin{tabular}{lcccc}
 & $(E, 0)$ & $(E, 1)$ & $(\sigma_{y}, 0)$ & $(\sigma_{y}, 1)$\\
$(E, 1)$ & $(E, 0)$ & $(E, 1)$ & $(\sigma_{y}, 1)$ & $(\sigma_{y}, 0)$\\
$(\sigma_{y}, 0)$ & $(\sigma_{y}, 1)$ & $(\sigma_{y}, 0)$ & $(E, 1)$ & $(E, 0)$\\
$(\sigma_{y}, 1)$ & $(\sigma_{y}, 0)$ & $(\sigma_{y}, 1)$ & $(E, 0)$ & $(E, 1)$
\end{tabular}
\end{center}
由此看出 $\overline{\mathbf{G}}^{B*}$ 同构于 $C_{4}$。我们关心的是 $\overline{\mathbf{G}}^{B*}$ 中满足 $\mathbf{D}(E, 1) = -\mathbf{1}$ 的那些表示（见式 (3.7.33)，取 $\alpha = 1$，$g = 2$）。由表 2.2 可见这是两个一维表示：
\begin{center}
\begin{tabular}{lcccc}
 & $(E, 0)$ & $(\sigma_{y}, 0)$ & $(E, 1)$ & $(\sigma_{y}, 1)$\\
${}^{1}E$ & $1$ & $-i$ & $-1$ & $i$\\
${}^{2}E$ & $1$ & $i$ & $-1$ & $-i$
\end{tabular}
\end{center}
由式 (3.8.8) 与 (3.8.9) 可见，$\mathbf{G}^{B}$ 的小表示的特征标表中相应的项为
\begin{center}
\begin{tabular}{lcc}
$B$ & $\{E \mid 000\}$ & $\{\sigma_{y} \mid \frac{1}{4}\frac{1}{4}\frac{1}{4}\}$\\
${}^{1}E$ & $1$ & $-i\zeta$\\
${}^{2}E$ & $1$ & $i\zeta$
\end{tabular}
\end{center}
其中 $\zeta = \exp(-\mathrm{i}\mathbf{k}_{1} \cdot \mathbf{v})$。

\subsection*{Q}

$\mathbf{G}^{Q}$ 的陪集代表：$\{E \mid 000\}$，$\{C_{2f} \mid {\frac{1}{4}}{\frac{1}{4}}{\frac{1}{4}}\}$。

$EQ = Q$，于是 $\mathbf{g}_{E} = 0$。

$C_{2f}Q = Q - \mathbf{g}_{1} - \mathbf{g}_{2} - \mathbf{g}_{3}$，于是 $\mathbf{g}_{Q} = -(\mathbf{g}_{1} + \mathbf{g}_{2} + \mathbf{g}_{3}) = (-1, -1, -1)$。

因子系统：$\mu(E, E) = 1$，$\mu(E, C_{2f}) = 1$，$\mu(C_{2f}, E) = 1$，以及 $\mu(C_{2f}, C_{2f}) = -i$。

中央延拓群 $\overline{\mathbf{G}}^{Q*}$ 由八个元素组成，同构于 $C_{8}$，使得 $(C_{2f}, 0)^{8} = (E, 0)$ 且 $(C_{2f}, 0)^{6} = (E, 1)$。我们关心的是 $\overline{\mathbf{G}}^{Q*}$ 中满足 $\mathbf{D}(E, 1) = i\mathbf{1}$ 的表示（见式 (3.7.33)，取 $\alpha = 1$，$g = 4$）。这是两个一维表示，其中特征标表的相应项为
\begin{center}
\begin{tabular}{lcc}
 & $(E, 0)$ & $(C_{2f}, 0)$\\
${}^{1}E_{3}$ & $1$ & $\exp(3\pi i/4)$\\
${}^{2}E_{1}$ & $1$ & $\exp(7\pi i/4)$
\end{tabular}
\end{center}
由于 $\exp(-\mathrm{i}\mathbf{k}_{1} \cdot \mathbf{v}) = \exp(-3\pi\mathrm{i}/4)$，可见 $\mathbf{G}^{Q}$ 的小表示的特征标表中相应的项为
\begin{center}
\begin{tabular}{lcc}
$Q$ & $\{E \mid 000\}$ & $\{C_{2f} \mid {\frac{1}{4}}{\frac{1}{4}}{\frac{1}{4}}\}$\\
${}^{1}E_{3}$ & $1$ & $+1$\\
${}^{2}E_{1}$ & $1$ & $-1$
\end{tabular}
\end{center}

\subsection*{Z}

$\mathbf{G}^{Z}$ 的陪集代表：$\{E \mid 000\}$，$\{C_{2x} \mid 000\}$，$\{\sigma_{y} \mid {\frac{1}{4}}{\frac{1}{4}}{\frac{1}{4}}\}$，$\{\sigma_{z} \mid {\frac{1}{4}}{\frac{1}{4}}{\frac{1}{4}}\}$。

\begin{center}
$\mathbf{g}_{E} = 0$，$\mathbf{g}_{C_{2x}} = (-1, 0, -1)$，$\mathbf{g}_{\sigma_{z}} = 0$，$\mathbf{g}_{\sigma_{y}} = (-1, 0, -1)$。
\end{center}

因子系统的元素全部等于 1，除了 $\mu(C_{2x}, \sigma_{z}) = \mu(C_{2x}, \sigma_{y}) = \mu(\sigma_{y}, \sigma_{z}) = \mu(\sigma_{y}, \sigma_{y}) = -1$。

中央延拓群 $\overline{\mathbf{G}}^{Z*}$ 由八个元素组成，同构于 $422\ (D_{4})$，使得 $(\sigma_{y}, 0)^{4} = (\sigma_{z}, 0)^{2} = (E, 0)$ 且 $(\sigma_{z}, 0)(\sigma_{y}, 0) = (\sigma_{y}, 0)^{3}$，$(\sigma_{z}, 0)(\sigma_{y}, 0)^{-1} = (C_{2x}, 0)$。我们关心的是 $\overline{\mathbf{G}}^{Z*}$ 中满足 $\mathbf{D}(E, 1) = -\mathbf{1}$ 的表示；由表 2.2 可见这是表示 $E$，其特征标为
\begin{center}
\begin{tabular}{lccccc}
 & $(E, 0)$ & $(E, 1)$ & $(\sigma_{y}, 0)(\sigma_{y}, 1)$ & $(\sigma_{z}, 0)(\sigma_{z}, 1)$ & $(C_{2x}, 0)(C_{2x}, 1)$\\
$E$ & $2$ & $-2$ & $0$ & $0$ & $0$
\end{tabular}
\end{center}
可能的矩阵为
\begin{equation*}
\mathbf{D}_{E}(E, 0) = \begin{pmatrix} 1 & 0\\ 0 & 1 \end{pmatrix}, \qquad
\mathbf{D}_{E}(\sigma_{y}, 0) = \begin{pmatrix} 0 & 1\\ -1 & 0 \end{pmatrix},
\end{equation*}
\begin{equation*}
\mathbf{D}_{E}(\sigma_{z}, 0) = \begin{pmatrix} 1 & 0\\ 0 & -1 \end{pmatrix}, \qquad
\mathbf{D}_{E}(C_{2x}, 0) = \begin{pmatrix} 0 & 1\\ 1 & 0 \end{pmatrix}.
\end{equation*}
记 $\zeta = \exp(-\mathrm{i}\mathbf{k}_{1} \cdot \mathbf{v})$，得到 $\overline{\mathbf{G}}^{Z}$ 的小表示的矩阵表示表中相应的项：
\begin{center}
\begin{tabular}{lcccc}
$Z$ & $\{E \mid 000\}$ & $\{\sigma_{y} \mid {\frac{1}{4}}{\frac{1}{4}}{\frac{1}{4}}\}$ & $\{\sigma_{z} \mid {\frac{1}{4}}{\frac{1}{4}}{\frac{1}{4}}\}$ & $\{C_{2x} \mid 000\}$\\[0.5ex]
$E$ & $\begin{pmatrix} 1 & 0\\ 0 & 1 \end{pmatrix}$ &
$\begin{pmatrix} 0 & \zeta\\ -\zeta & 0 \end{pmatrix}$ &
$\begin{pmatrix} \zeta & 0\\ 0 & -\zeta \end{pmatrix}$ &
$\begin{pmatrix} 0 & 1\\ 1 & 0 \end{pmatrix}$
\end{tabular}
\end{center}

\subsection*{S}

$\mathbf{G}^{S}$ 的陪集代表：$\{E \mid 000\}$，$\{C_{2c} \mid {\frac{1}{4}}{\frac{1}{4}}{\frac{1}{4}}\}$，$\{\sigma_{de} \mid 000\}$，$\{\sigma_{y} \mid {\frac{1}{4}}{\frac{1}{4}}{\frac{1}{4}}\}$。

\begin{center}
$\mathbf{g}_{E} = 0$，$\mathbf{g}_{C_{2c}} = (-1, 0, -1)$，$\mathbf{g}_{\sigma_{de}} = 0$，$\mathbf{g}_{\sigma_{y}} = (-1, 0, -1)$。
\end{center}

因子系统的元素全部等于 1，除了 $\mu(C_{2c}, C_{2c}) = \mu(C_{2c}, \sigma_{y}) = \mu(\sigma_{y}, C_{2c}) = \mu(\sigma_{y}, \sigma_{y}) = -1$。

中央延拓群 $\overline{\mathbf{G}}^{S*}$ 由八个元素组成，同构于 $4/m\ (C_{4h})$，使得 $(\sigma_{y}, 0)^{4} = (\sigma_{de}, 0)^{2} = (E, 0)$ 且 $(\sigma_{y}, 0)(\sigma_{de}, 0) = (\sigma_{de}, 0)(\sigma_{y}, 0) = (C_{2c}, 0)$。我们关心的是 $\overline{\mathbf{G}}^{S*}$ 中满足 $\mathbf{D}(E, 1) = -\mathbf{1}$ 的那些表示；由表 2.2 可见这是四个一维表示，其相应特征标为
\begin{center}
\begin{tabular}{lcccc}
 & $(E, 0)$ & $(\sigma_{de}, 0)$ & $(\sigma_{y}, 0)$ & $(C_{2c}, 0)$\\
${}^{1}E_{1}$ & $1$ & $1$ & $-i$ & $-i$\\
${}^{2}E_{1}$ & $1$ & $1$ & $i$ & $i$\\
${}^{2}E_{2}$ & $1$ & $-1$ & $i$ & $-i$\\
${}^{1}E_{2}$ & $1$ & $-1$ & $-i$ & $i$
\end{tabular}
\end{center}
记 $\zeta = \exp(-\mathrm{i}\mathbf{k}_{1} \cdot \mathbf{v})$，得到 $\mathbf{G}^{S}$ 的小表示的特征标表中相应的项：
\begin{center}
\begin{tabular}{lcccc}
$S$ & $\{E \mid 000\}$ & $\{\sigma_{de} \mid 000\}$ & $\{\sigma_{y} \mid {\frac{1}{4}}{\frac{1}{4}}{\frac{1}{4}}\}$ & $\{C_{2c} \mid {\frac{1}{4}}{\frac{1}{4}}{\frac{1}{4}}\}$\\
${}^{1}E_{1}$ & $1$ & $1$ & $-i\zeta$ & $i\zeta$\\
${}^{2}E_{1}$ & $1$ & $1$ & $i\zeta$ & $i\zeta$\\
${}^{2}E_{2}$ & $1$ & $-1$ & $i\zeta$ & $-i\zeta$\\
${}^{1}E_{2}$ & $1$ & $-1$ & $-i\zeta$ & $i\zeta$
\end{tabular}
\end{center}

我们用来求非点式空间群布里渊区面上波矢 $\mathbf{k}$ 的小表示的方法，当然也可用于对称点。然而有一种 Herring (1942) 的替代方法，它避免了对这类对称点使用射影表示；由于其做法更为直接（它使用小群的因子群而不是小陪群），我们将描述它，并以它为例应用于 $Fd3m\ (O^{7}_{h})$ 的点 $W$、$X$ 和 $L$。两种方法难分高下，因为它们都要求一个高阶群的特征标。

Herring 的方法如下。在任何小表示中
\begin{equation*}
\Gamma^{\mathbf{k}_{1}}_{p}(\{E \mid \mathbf{t}\}) = \exp(-\mathrm{i}\mathbf{k}_{1} \cdot \mathbf{t})\mathbf{1}.
\tag{3.8.10}
\end{equation*}
若 $\mathbf{k}_{1}$ 是一个对称点，则对大量的平移 $\mathbf{t}$ 有相位因子 $\exp(-\mathrm{i}\mathbf{k}_{1} \cdot \mathbf{t}) = 1$。把 $\mathbf{T}$ 中具有这一性质的元素构成的子群记为 $\mathbf{T}^{\mathbf{k}_{1}}$。则 $\mathbf{T}^{\mathbf{k}_{1}}$ 是 $\mathbf{G}^{\mathbf{k}_{1}}$ 的不变子群，而因子群 $\mathbf{G}^{\mathbf{k}_{1}}/\mathbf{T}^{\mathbf{k}_{1}}$（记作 ${}^{H}\mathbf{G}^{\mathbf{k}_{1}}$）的阶比 $\mathbf{G}^{\mathbf{k}_{1}}$ 的阶小（在确定空间群单值表示时出现的 ${}^{H}\mathbf{G}^{\mathbf{k}_{1}}$ 的最大可能阶是 96）。现在由于对一切 $\mathbf{t} \in \mathbf{T}^{\mathbf{k}_{1}}$ 有 $\Gamma^{\mathbf{k}_{1}}_{p}(\{E \mid \mathbf{t}\}) = \mathbf{1}$，可知 $\mathbf{G}^{\mathbf{k}_{1}}$ 中同一 $\mathbf{T}^{\mathbf{k}_{1}}$ 陪集内的每个元素都被同一矩阵表示。于是，由于 $\mathbf{T}^{\mathbf{k}_{1}}$ 的陪集就是 ${}^{H}\mathbf{G}^{\mathbf{k}_{1}}$ 的元素，$\Gamma^{\mathbf{k}_{1}}_{p}$ 可以拿来作为该因子群的表示。粗略地说，我们把同一陪集的所有元素都等同起来。反之，利用这种等同，我们也可以把因子群的表示当作小群的表示。当然，由于式 (3.8.10) 施加的额外要求，其中只有一部分是小表示；但所有小表示肯定都可以这样通过把它们从 ${}^{H}\mathbf{G}^{\mathbf{k}_{1}}$ 提升到 $\mathbf{G}^{\mathbf{k}_{1}}$ 而得到，因为小表示总具有性质：对一切 $\mathbf{t} \in \mathbf{T}^{\mathbf{k}_{1}}$ 有 $\Gamma^{\mathbf{k}_{1}}_{p}(\{E \mid \mathbf{t}\}) = \mathbf{1}$。

为了使工作显得自然，习惯上用陪集代表作为陪集的符号。这些符号随后构成一个同构于所给因子群的群。使用它们时，我们把仅相差 $\mathbf{t} \in \mathbf{T}^{\mathbf{k}_{1}}$ 型平移的任意两个符号视为相同。由于这可以一目了然地做到，因子群内的乘法可以立即由普通的空间群乘法规则推得。

回到金刚石的情形，现在说明 Herring 方法对点 $W$、$X$ 和 $L$ 如何运作。

\subsection*{W}

由于 $W = \frac{1}{2}\mathbf{g}_{1} + \frac{1}{4}\mathbf{g}_{2} + \frac{3}{4}\mathbf{g}_{3}$，群 ${}^{H}\mathbf{G}^{W}$ 含有四个平移 $\{E \mid \mathbf{0}\}$、$\{E \mid \mathbf{t}_{1}\}$、$\{E \mid \mathbf{t}_{2}\}$、$\{E \mid \mathbf{t}_{3}\}$，分别由 $\mathbf{1}$、$-\mathbf{1}$、$-i\mathbf{1}$ 和 $i\mathbf{1}$ 表示。于是 ${}^{H}\mathbf{G}^{W}$ 中有 32 个元素，即八个元素 $\{E \mid \mathbf{0}\}$、$\{C_{2x} \mid \mathbf{0}\}$、$\{C_{2a} \mid \mathbf{v}\}$、$\{C_{2f} \mid \mathbf{v}\}$、$\{C_{2f} \mid \bar{\mathbf{v}}\}$、$\{\sigma_{y} \mid \mathbf{v}\}$、$\{\sigma_{z} \mid \mathbf{v}\}$、$\{S^{+}_{4x} \mid \mathbf{0}\}$、$\{S^{-}_{4x} \mid \mathbf{0}\}$ 与这四个平移在 Herring 意义下的乘积。这个 32 阶的群可以辨认为表 5.1 中的 $G^{4}_{32}$，取 $\{S^{+}_{4x} \mid \mathbf{0}\} = P$、$\{C_{2f} \mid \mathbf{v}\} = R$、$\{E \mid \mathbf{t}_{3}\} = Q$。（只需检验生成关系 $P^{4} = R^{2} = Q^{4} = E$、$PQ = QP$、$RQ = QR$ 以及 $RP = Q^{3}P^{3}R$ 即可。）我们需要的是 $\{E \mid 001\} = Q$ 由 $i\mathbf{1}$ 表示的那些表示。这是两个二维表示。${}^{H}\mathbf{G}^{W}$ 的特征标表的有关部分为

\begin{figure}[H]
\centering
\includegraphics[width=0.92\textwidth]{c3_wchar_p1.png}
\end{figure}
\begin{figure}[H]
\centering
\includegraphics[width=0.92\textwidth]{c3_wchar_p2.png}
\caption*{\textbf{W 点}\ ${}^{H}\mathbf{G}^{W}$ 的特征标表（原书扫描占位，跨第 168--169 页）。}
\end{figure}

生成这些表示的合适矩阵是
\begin{center}
\begin{tabular}{lcc}
 & $\{C_{2f} \mid {\frac{1}{4}}{\frac{1}{4}}{\frac{1}{4}}\}$ & $\{S^{+}_{4x} \mid 000\}$\\[0.5ex]
${}^{1}\Gamma_{1}$ & $\begin{pmatrix} 0 & 1\\ 1 & 0 \end{pmatrix}$ & $\begin{pmatrix} 1 & 0\\ 0 & -i \end{pmatrix}$\\[2ex]
${}^{1}\Gamma_{2}$ & $\begin{pmatrix} 0 & 1\\ 1 & 0 \end{pmatrix}$ & $\begin{pmatrix} i & 0\\ 0 & -1 \end{pmatrix}$
\end{tabular}
\end{center}

\subsection*{X}

由于 $X = (\frac{1}{2}0\frac{1}{2})$，群 ${}^{H}\mathbf{G}^{X}$ 含有两个平移 $\{E \mid \mathbf{0}\}$ 与 $\{E \mid \mathbf{t}_{1}\}$，分别由 $\mathbf{1}$ 与 $-\mathbf{1}$ 表示。于是 ${}^{H}\mathbf{G}^{X}$ 中有 32 个元素，即 16 个元素 $\{E \mid \mathbf{0}\}$、$\{C_{2y} \mid \mathbf{0}\}$、$\{C^{+}_{4y} \mid \mathbf{v}\}$、$\{C_{2x} \mid \mathbf{0}\}$、$\{C_{2z} \mid \mathbf{0}\}$、$\{C_{2c} \mid \mathbf{v}\}$、$\{C_{2e} \mid \mathbf{v}\}$、$\{I \mid \mathbf{v}\}$、$\{\sigma_{y} \mid \mathbf{v}\}$、$\{S^{+}_{4y} \mid \mathbf{0}\}$、$\{\sigma_{x} \mid \mathbf{v}\}$、$\{\sigma_{z} \mid \mathbf{v}\}$、$\{\sigma_{dc} \mid \mathbf{0}\}$ 与 $\{\sigma_{de} \mid \mathbf{0}\}$ 与这两个平移在 Herring 意义下的乘积。这个 32 阶的群可以辨认为表 5.1 中的 $G^{2}_{32}$，取 $\{\sigma_{x} \mid \mathbf{v}\} = P$、$\{S^{+}_{4y} \mid \mathbf{0}\} = Q$、$\{C_{2x} \mid \mathbf{0}\} = R$。我们需要的是 $\{E \mid \mathbf{t}_{1}\} = P^{2}$ 由 $-\mathbf{1}$ 表示的那些表示。这是四个二维表示。${}^{H}\mathbf{G}^{X}$ 的特征标表的有关部分为

\begin{figure}[H]
\centering
\includegraphics[width=0.92\textwidth]{c3_xchar.png}
\caption*{\textbf{X 点}\ ${}^{H}\mathbf{G}^{X}$ 的特征标表（原书扫描占位）。}
\end{figure}

生成这些表示的合适矩阵是
\begin{center}
\begin{tabular}{lccc}
 & $\{\sigma_{x} \mid {\frac{1}{4}}{\frac{1}{4}}{\frac{1}{4}}\}$ & $\{S^{+}_{4y} \mid 000\}$ & $\{C_{2x} \mid 000\}$\\[0.5ex]
$E_{1}$ & $\begin{pmatrix} 0 & -1\\ 1 & 0 \end{pmatrix}$ & $\begin{pmatrix} 0 & 1\\ 1 & 0 \end{pmatrix}$ & $\begin{pmatrix} 0 & 1\\ 1 & 0 \end{pmatrix}$\\[2ex]
$E_{2}$ & $\begin{pmatrix} 0 & -1\\ 1 & 0 \end{pmatrix}$ & $\begin{pmatrix} 0 & 1\\ -1 & 0 \end{pmatrix}$ & $\begin{pmatrix} 0 & 1\\ 1 & 0 \end{pmatrix}$\\[2ex]
$E_{3}$ & $\begin{pmatrix} 0 & -1\\ 1 & 0 \end{pmatrix}$ & $\begin{pmatrix} 0 & -1\\ -1 & 0 \end{pmatrix}$ & $\begin{pmatrix} 0 & 1\\ 1 & 0 \end{pmatrix}$\\[2ex]
$E_{4}$ & $\begin{pmatrix} 0 & -1\\ 1 & 0 \end{pmatrix}$ & $\begin{pmatrix} 0 & -1\\ 1 & 0 \end{pmatrix}$ & $\begin{pmatrix} 0 & 1\\ 1 & 0 \end{pmatrix}$
\end{tabular}
\end{center}

\subsection*{L}

由于 $L = (\frac{1}{2}\frac{1}{2}\frac{1}{2})$，群 ${}^{H}\mathbf{G}^{L}$ 含有两个平移 $\{E \mid \mathbf{0}\}$ 与 $\{E \mid \mathbf{t}_{1}\}$，分别由 $\mathbf{1}$ 与 $-\mathbf{1}$ 表示。该群阶为 24，同构于 $3m\ (C_{3v}) \otimes \mathbf{J} \otimes \mathbf{T}_{1}$，其中 $3m\ (C_{3v})$ 含有六个元素 $\{E \mid \mathbf{0}\}$、$\{C^{\pm}_{31} \mid \mathbf{0}\}$、$\{\sigma_{db} \mid \mathbf{0}\}$、$\{\sigma_{dc} \mid \mathbf{0}\}$、$\{\sigma_{df} \mid \mathbf{0}\}$，$\mathbf{J}$ 含有 $\{E \mid \mathbf{0}\}$ 与 $\{I \mid \mathbf{v}\}$，$\mathbf{T}_{1}$ 含有两个平移 $\{E \mid \mathbf{0}\}$ 与 $\{E \mid \mathbf{t}_{1}\}$。因此有六个小表示，${}^{H}\mathbf{G}^{L}$ 的特征标表的有关部分由表 2.2 可见为
\begin{center}
\begin{tabular}{lcccccc}
$L$ & $\{E \mid 000\}$ & $\{C^{\pm}_{31} \mid 000\}$ & $\begin{array}{c}\{\sigma_{db} \mid 000\}\\ \{\sigma_{dc} \mid 000\}\\ \{\sigma_{df} \mid 000\}\end{array}$ & $\{I \mid {\frac{1}{4}}{\frac{1}{4}}{\frac{1}{4}}\}$ & $\{S^{\mp}_{61} \mid {\frac{1}{4}}{\frac{1}{4}}{\frac{1}{4}}\}$ & $\begin{array}{c}\{C_{2b} \mid {\frac{1}{4}}{\frac{1}{4}}{\frac{1}{4}}\}\\ \{C_{2e} \mid {\frac{1}{4}}{\frac{1}{4}}{\frac{1}{4}}\}\\ \{C_{2f} \mid {\frac{1}{4}}{\frac{1}{4}}{\frac{1}{4}}\}\end{array}$\\
$A_{1g}$ & $1$ & $1$ & $1$ & $1$ & $1$ & $1$\\
$A_{2g}$ & $1$ & $1$ & $-1$ & $1$ & $1$ & $-1$\\
$E_{g}$ & $2$ & $-1$ & $0$ & $2$ & $-1$ & $0$\\
$A_{1u}$ & $1$ & $1$ & $1$ & $-1$ & $-1$ & $-1$\\
$A_{2u}$ & $1$ & $1$ & $-1$ & $-1$ & $-1$ & $1$\\
$E_{u}$ & $2$ & $-1$ & $0$ & $-2$ & $1$ & $0$
\end{tabular}
\end{center}
二维表示的生成矩阵可由表 2.3 得到。可以看到，碰巧 $Fd3m$ 中属于 $Q$ 和 $L$ 的表示的形式与 $Fm3m$ 中的完全一样。我们给出的关于这种情况何时必然发生的规则因此是充分的但不必要的；从这个例子我们看到，即使 $\mathbf{k}_{1}$ 位于布里渊区面上且空间群是非点式的，$\mathbf{G}^{\mathbf{k}_{1}}$ 的小表示偶尔也与 $\overline{\mathbf{G}}^{\mathbf{k}_{1}}$ 的表示具有相同的形式。
